\scanpage{43}
\mspage{17}

En fait, les antiinvolutions de $\mathbb C$ ont toutes des pts fixes, et ce sont les symétries orth. p.r. aux droites réelles (pas néc. passant par l'origine) \deleted{\unsure{par}}dans $\mathbb C$. La surface à bord conforme correspondante est isom. à un \uuline{demi-plan} (avec bord). On note que tte métrique \inserted{sur $\mathscr X$} quasi-invariante par autom. conforme, est invariante par autom d'ordre fini, en particulier par $\sigma$. Donc il y a sur $X$ une classe \inserted{can} de métriques riemanniennes, mod. cte.

\uline{NB} \deleted{\illegible{1}} Le modèle des $X$ envisagés est le demi-plan supérieur $\deleted{\unsure{\mathbb C|G}}$
\[
\overline{\mathbb H}=\{z\in\mathbb C\mid\unsure{\mathfrak I|\mathfrak J|\mathcal J}z\geq0\}.
\]
Le groupe des autom de $\overline{\mathbb H}$, isom au groupe des autom de la droite affine $\mathbb C$ qui commutent à $\sigma:z\longmapsto\bar z$, est \uline{le groupe affine réel}
% archive sheet number: 45
